%!TEX root = ../graduate-work.tex

\section{Постановка задачи}

% Здесь надо максимально формально описать суть задачи, которую потребуется решить, так, чтобы можно было потом понять, в какой степени полученное в результате работы решение ей соответствует. Текст главы должен быть написан в стиле технического задания, т.е. содержать как описание задачи, так и некоторый набор требований к решению

% \textbf{Алгоритм Решения Задачи.}
% \begin{itemize}
% \item[1.] Описать Задачу
% \item[2.] Декомпозировать задачу $a=mH\in\mathcal
% C,\;wt(e')=\lfloor\frac{d-1}2\rfloor;$
% \item[3.] Решить каждую подзадачу и объединить решение.
% \end{itemize}

Задачей выпускной квалификационной работы является разработка метода оценки качества русскоязычных радиологических отчётов органов грудной клетки в условиях ограниченного объёма размеченных данных.

Особенность рассматриваемой задачи состоит в том, что для русскоязычной медицинской предметной области объём доступных данных слишком мал для обучения полноценной специализированной модели оценки качества. В связи с этим в работе предлагается не обучение новой модели с нуля, а построение метода, основанного на комбинировании существующих текстовых, визуально-текстовых и клинически-ориентированных методов оценки с автоматическим подбором их конфигурации и весов на обучающей выборке.

В рамках задачи можно выделить несколько подзадач:
\begin{enumerate}
    \item Изучение существующих методов оценки качества радиологических отчётов и выбор набора базовых подходов, пригодных для адаптации к использованию на русском языке.
    \item Формирование экспериментального набора данных для оценки качества русскоязычных радиологических отчётов.
    \item Разработка метода, комбинирующего базовые методы оценки, методику их выбора и смешивания результатов.
    \item Программная реализация предложенного подхода.
    \item Анализ полученных результатов, сравнение с отдельными базовыми способами оценки и интерпретация качества итогового метода с точки зрения согласованности с врачебной оценкой.
\end{enumerate}